\scanpage{65}
6 ps.dr. qui joignent 2 à 2 les quatre sommets
$u,u',r,t$ de la figure. Mais on constate alors que
la bijection \inserted{$\varphi_s:$} $\Omega(D)\simeq\Omega(D')$ définie par $s$ n'est autre
que celle définie par $r$, qui est dans la \unsure{m\textsuperscript{e}|$=$}
des deux régions de $X$ découpées par $D,D'$\ldots)

\begin{center}\rule{1cm}{.3pt}\end{center}

$K^1$, 1 complexe subdivisé en ps. droites (\unsure{tout|tt} sommet
appartient à au moins 2 ps.dr, et 2 ps.droites se rencontrent en un pt et un seul)
Si $D,D'$ sont 2 ps.dr. distinctes, et $s$ un sommet
qui n'est ni sur $D$, ni sur $D'$, on a une bijection
\[
\varphi_{s;\deleted{\unsure{g|D}}D,D'}\ \Omega(D)\simeq\Omega(D')
\]
avec les propriétés suivantes :

a) Sur l'ens des ps. droites \inserted{$D,D'\ldots$} qui ne passent pas
par $s$, \deleted{\unsure{l}} le syst. des $\varphi_{s;D,D'}$ est simplement transitif.

b) Si $\deleted{\unsure{D''}}\Delta$ est une ps. droite distincte de $D$ et $D'$,
coup $D$ et $D'$ en $x$ et $x'$, \inserted{(\uline{NB} on n'exclut pas $x=x'$)}, et si $s,s'$ sont deux
sommets \deleted{sur} $\Delta$ distincts de $x,x'$, \deleted{\unsure{et sur|alors}} \inserted{alors}
\deleted{un m\textsuperscript{e} des} on a $\varphi_{s;D,D'}=\varphi_{s';D,D'}$ ss $s,s'$
appartiennent au m\textsuperscript{e} arc des deux arcs définis par :
$x$ et $x'$ sur $\Delta$. Si $a,a'$ sont \deleted{les} deux arcs,
on peut donc définir $\varphi_{a;D,D'}$ et $\varphi_{a';D,D'}$

c) Soit $t$ \deleted{\unsure{le sommet}} l'intersection de $D$ et $D'$.
et soient $\Delta_i$ \deleted{$\Delta_2$} \inserted{$(i=1,2)$} deux psdr. passant par $s$ et non par
$t$, \deleted{\unsure{soi}} $u_i$ (resp $u'_i$) $(i=1,2)$ \unsure{ses|son|leurs} int. avec $D$ $(D')$.
% archive sheet number: 64
